% Adapted from week6_lecture2.tex; inherited harmonic and periodicity errors corrected.

\section{Introduction and Review}
We start by reviewing concepts we have discussed in class before: (1) periodicity of signals, (2) properties of convolution, and (3) the eigenfunction property of complex exponentials for LTI systems. 
\subsection{Periodicity of signals}
Recall the periodicity of signals by solving the popquiz below. Assume $\omega_0>0$ for all problems.
\begin{popquiz}
  Prove that $x(t)=e^{j\omega t}$ is a periodic signal. Find its fundamental period.
  \popqsplit 
  Periodicity requires us to prove that there exists a $T>0$ such that $x(t+T)=x(t)$ for all $t$. We have
  \[
  x(t+T)=e^{j\omega (t+T)}=e^{j\omega t}e^{j\omega T}=x(t)e^{j\omega T}.
  \]
  Therefore, we need to find a $T>0$ such that $e^{j\omega T}=1$. This is true if and only if $|\omega|T=2\pi k$ for some integer $k>0$.
  Therefore, we can choose $T=2\pi k/|\omega|$. The fundamental period is the smallest positive $T$, which occurs when $k=1$, so the fundamental period is $T=2\pi/|\omega|$.
\end{popquiz}
As a follow-up, solve another popquiz! Once you're done, you can review the solutions at the end.
\begin{popquiz}
Prove that the following signals are period:
\begin{enumerate}
\item Scaled harmonic of the complex exponential: $x(t)=e^{jk\omega_0 t}$, where $k$ is a non-zero integer.
\item Finite harmonic sum: $x(t)=\sum_{k=-K}^{K}a_k e^{jk\omega_0t}$, where $K$ is a positive integer and $k$ is a non-zero integer.
\end{enumerate}
\popqsplit
For $x(t)=e^{jk\omega_0 t}$, we have
\[
x(t+T)=e^{jk\omega_0 (t+T)}=e^{jk\omega_0 t}e^{jk\omega_0 T}=x(t)e^{jk\omega_0 T}.
\]
Therefore, we need to find a $T>0$ such that $e^{jk\omega_0 T}=1$. This is true if and only if $|k\omega_0|T=2\pi m$ for some integer $m>0$. Therefore, we can choose $T=2\pi m/|k\omega_0|$. The fundamental period is the smallest positive $T$, which occurs when $m=1$, so the fundamental period is $T=2\pi/|k\omega_0|$. The second signal (summed harmonics) is also periodic with period $T=2\pi/\omega_0$ since each term in the sum is periodic with that period, so the sum is also periodic with the same period.
\end{popquiz}
\subsection{Properties of convolution}
We note two important properties of convolution that often come in handy: commutativity and associativity. For two signals $x(t)$ and $h(t)$, we have
\[
x(t)*h(t)=h(t)*x(t),
\]
How can we prove this? Hint: Use the change of variable (colloquially known as ``u-substitution'').

To prove commutativity, we can start with the definition of convolution for the left-hand side:
\[
x(t)*h(t)=\int_{-\infty}^{\infty} x(\tau)h(t-\tau)\,d\tau.
\]
Now, substitute $u = t - \tau$. Then, $du = -d\tau$ because $t$ is a fixed value for the function inside the integral. The limits of integration become: $+\infty$ to $-\infty$ when $\tau$ goes from $-\infty$ to $+\infty$. Therefore, we have
\[
x(t)*h(t)=\int_{+\infty}^{-\infty} x(t-u)h(u)(-du)=\int_{-\infty}^{+\infty} h(u)x(t-u)\,du=h(t)*x(t).
\]

Associativity of convolution can be proved similarly. For three signals $x_1(t)$, $h(t)$, and $x_2(t)$, we have
\[
x_1(t)*(h(t)*x_2(t))=(x_1(t)*h(t))*x_2(t).
\]
To prove associativity, we can start with the left-hand side:
\[
x_1(t)*(h(t)*x_2(t))=\int_{-\infty}^{\infty} x_1(\tau)\left[\int_{-\infty}^{\infty} h(\lambda)x_2(t-\tau-\lambda)\,d\lambda\right]\,d\tau.
\]
Now, we can change the order of integration:
\[
x_1(t)*(h(t)*x_2(t))=\int_{-\infty}^{\infty} h(\lambda)\left[\int_{-\infty}^{\infty} x_1(\tau)x_2(t-\tau-\lambda)\,d\tau\right]\,d\lambda.
\] 
The inner integral is the convolution of $x_1(t)$ and $x_2(t)$ evaluated at $t-\lambda$. Therefore, we have
\[
x_1(t)*(h(t)*x_2(t))=\int_{-\infty}^{\infty} h(\lambda)(x_1*x_2)(t-\lambda)\,d\lambda=(x_1*x_2)*h(t).
\]

\subsection{Eigenfunctions of LTI systems}
At the end of the previous lecture, we established that $e^{j\omega t}$ is an eigenfunction of LTI systems. That is, for an input $x(t)=e^{j\omega t}$, the output $y(t)$ is also a complex exponential at the same frequency $\omega$ but scaled by a complex number $H(j\omega)$:
\[
y(t)=H(j\omega)e^{j\omega t}.
\]
where
\[
H(j\omega)=\int_{-\infty}^{\infty} h(\tau)\,e^{-j\omega \tau}\,d\tau.
\]
As a result of this \emph{very important} result, we proposed that if we are able to write any signal $x(t)$ as a linear combination of complex exponentials, then we can find the output $y(t)$ of an LTI system by simply scaling each complex exponential by $H(j\omega)$ and adding them up!

The previous line is a one-line summary of Fourier analysis and synthesis --- something that we will be spending a lot of time on in the next few weeks.

We start this journey by positing the following problem:  Given a $T$-periodic signal $x(t)$, can we write it as a linear combination of complex exponentials? If so, how? More concretely, can we write
\[
x(t)=\sum_{k=-\infty}^{\infty} a_k\,e^{jk\omega_0 t},
\qquad \omega_0=\tfrac{2\pi}{T},
\]
the complex Fourier series coefficients are $\{a_k\}$. Note that $a_k \in \mathbb{C}$ are complex numbers, in general. The big question is; how do we find $a_k$?

\section{Goals}
Represent a periodic signal $x(t)$ as a linear combination of harmonically related complex exponentials:
\begin{equation}
x(t)=\sum_{k=-\infty}^{\infty} a_k\,e^{jk\omega_0 t}.
\label{eq:fourier-series}
\end{equation}
Find the coefficients $a_k$.
\section{Introduction to Fourier Series}
Since the first part of the goal is ``any periodic signal $x(t)$'' and we are claiming that $x(t)$ can be written as a linear combination of complex exponentials. So, it must be true that (and we should make sure of it) that $e^{jk\omega_0 t}$ is periodic for every integer $k$.
% \begin{popquiz}
% Let $\omega_0>0$ and $T=2\pi/\omega_0$.
% \begin{enumerate}
%   \item Find the fundamental period of $e^{j\omega_0t}$ and of
%   $e^{jk\omega_0t}$ for an integer $k\ne0$. What happens when $k=0$?
%   \item Show that a finite harmonic sum
%   $\sum_{k=-K}^{K}a_k e^{jk\omega_0t}$ has $T$ as a period.
%   Must $T$ be its fundamental period?
% \end{enumerate}
% \popqsplit
% For $k\ne0$, periodicity requires $e^{jk\omega_0T_k}=1$, so the smallest
% positive period is $T_k=2\pi/(|k|\omega_0)=T/|k|$.
% For $k=0$, the signal is the constant $1$: every positive number is a
% period, and there is no smallest positive fundamental period.

% Each harmonic satisfies $e^{jk\omega_0(t+T)}=e^{jk\omega_0t}$ because
% $e^{j2\pi k}=1$. Therefore every finite harmonic sum has $T$ as a common
% period. Its fundamental period can be shorter: for example,
% $\cos(2\omega_0t)$ has fundamental period $T/2$.
% A common period should not automatically be called the fundamental period.
% \end{popquiz}

Let's start to answer the second part of the goal: how do we find $a_k$? We will start with a simple example.
\subsection{Example: A mix-sinusoidal audio signal}
Use an idealized combination of three tones to see how Fourier coefficients encode frequency, amplitude, and phase. Real musical notes usually contain several harmonics; the example below is a simple sum of sinusoids.

Consider
\[
x(t)=\sin(6t)+\cos(2t)+\sin(12t),\qquad \omega_0=2.
\]
Write $x(t)$ as a linear combination of $e^{jk\omega_0 t}$:
\[
\sin(6t)=\frac{1}{2j}\!\left(e^{j6t}-e^{-j6t}\right)
= \frac{1}{2j}\!\left(e^{j(3)\omega_0 t}-e^{-j(3)\omega_0 t}\right),
\]
\[
\cos(2t)=\frac{1}{2}\!\left(e^{j2t}+e^{-j2t}\right)
=\frac{1}{2}\!\left(e^{j(1)\omega_0 t}+e^{-j(1)\omega_0 t}\right),
\]
\[
\sin(12t)=\frac{1}{2j}\!\left(e^{j12t}-e^{-j12t}\right)
= \frac{1}{2j}\!\left(e^{j(6)\omega_0 t}-e^{-j(6)\omega_0 t}\right).
\]
Hence
\[
x(t)=\sum_{k=-6}^{6} a_k\,e^{jk\omega_0 t},\qquad \omega_0=2,
\]
with the nonzero Fourier series coefficients as
\[
\begin{aligned}
a_{\pm 1}=\frac{1}{2},\qquad a_{\pm 3}=\pm \frac{1}{2j},a_{\pm 6}=\pm \frac{1}{2j},
\end{aligned}
\]
where the ``$\pm$'' pairs obey $a_{-k}=a_k^\ast$ for this real $x(t)$.

The cosine contributes equal real coefficients at $k=\pm1$. The sine terms contribute opposite, purely imaginary coefficients at $k=\pm3,\pm6$. All other $a_k$ are zero. The fundamental period of this signal is $T=\pi$.

\section{The Fourier analysis equation}
\label{sec:w6l2-analysis}
To find $a_k$ generally, we have the Fourier series analysis equation (the equation that gives us values of $a_k$):
\[
\,a_n=\frac{1}{T}\int_{0}^{T} x(t)\,e^{-jn\omega_0 t}\,dt\,
\qquad n\in\mathbb{Z},
\]
if you replace $n$ by $k$, you get the more familiar form:
\[
a_k=\frac{1}{T}\int_{0}^{T} x(t)\,e^{-jk\omega_0 t}\,dt,
\qquad k\in\mathbb{Z},
\]
which is the formula for the Fourier series coefficients.
Note that for $k=0$, we have
\[
a_0=\frac{1}{T}\int_{0}^{T} x(t)\,dt,
\]
which is the average value (or DC component) of $x(t)$ over one period.

\section{Fourier series properties}
\label{sec:w6l2-properties}
There are many helpful properties that you should know about Fourier series coefficients. Here are a few of them (you can find the full list in the textbook):

\subsection{Linearity}

For two signals represented using the same period $T$ and harmonic frequencies, with Fourier series coefficients $a_k$ and $b_k$, if we construct another signal by linear superposition, $z(t)=A x(t)+B y(t)$, then the Fourier Series coefficients for $z(t)$ satisfy
\[
z(t)\;\Longleftrightarrow\; \{\,A a_k + B b_k\,\}_{k\in\mathbb{Z}}.
\]

We will talk about many more properties next week.

\section{Practice Problems}
\begin{enumerate}
  \item Review continuous-time Fourier series analysis, synthesis, and
  properties in Chapter 3 of Oppenheim and Willsky, \emph{Signals and Systems}
  (2nd edition), especially Sections 3.2--3.5 and Table 3.1.
  \item Practice finding Fourier coefficients by integrating over one
  period of a rectangular waveform and a periodic ramp.
  \item Use the LTI filtering rule to connect input and output Fourier
  coefficients, as required in Pset 6.
\end{enumerate}
\textbf{Reminder:} Pset 6 is due October 14, 2026.

For an additional reference on harmonic analysis and LTI filtering, see
\href{https://ocw.mit.edu/courses/6-003-signals-and-systems-fall-2011/72c3d01402fa6b7c109b7b88d399115e_MIT6_003F11_lec15.pdf}{MIT 6.003 Lecture 15: Fourier Series}.
